\documentclass{standalone}
\usepackage{circuitikz}
\usetikzlibrary{calc}
\definecolor{circuitnotemark}{HTML}{FE00FE}
\makeatletter
\pgfdeclareshape{dev4w11n8}{
  \anchor{center}{\pgfpointorigin}
  \anchor{text}{\pgfpoint{0.4cm-.5\wd\pgfnodeparttextbox}{-.5\ht\pgfnodeparttextbox}}
  \anchor{value}{\pgfpoint{0.4cm}{0cm}}
  \anchor{north}{\pgfpoint{0cm}{1cm}}
  \anchor{south}{\pgfpoint{0cm}{-1cm}}
  \anchor{east}{\pgfpoint{1.1cm}{0cm}}
  \anchor{west}{\pgfpoint{-1.1cm}{0cm}}
  \anchor{pin 1}{\pgfpoint{-1.5cm}{0.75cm}}
  \anchor{bpin 1}{\pgfpoint{-1.1cm}{0.75cm}}
  \anchor{pin 2}{\pgfpoint{-1.5cm}{0.25cm}}
  \anchor{bpin 2}{\pgfpoint{-1.1cm}{0.25cm}}
  \anchor{pin 3}{\pgfpoint{-1.5cm}{-0.25cm}}
  \anchor{bpin 3}{\pgfpoint{-1.1cm}{-0.25cm}}
  \anchor{pin 4}{\pgfpoint{-1.5cm}{-0.75cm}}
  \anchor{bpin 4}{\pgfpoint{-1.1cm}{-0.75cm}}
  \backgroundpath{
    \pgfpathrectanglecorners{\pgfpoint{-1.1cm}{-1cm}}{\pgfpoint{1.1cm}{1cm}}
    \pgfpathmoveto{\pgfpoint{-1.5cm}{0.75cm}}\pgfpathlineto{\pgfpoint{-1.1cm}{0.75cm}}
    \pgfpathmoveto{\pgfpoint{-1.5cm}{0.25cm}}\pgfpathlineto{\pgfpoint{-1.1cm}{0.25cm}}
    \pgfpathmoveto{\pgfpoint{-1.5cm}{-0.25cm}}\pgfpathlineto{\pgfpoint{-1.1cm}{-0.25cm}}
    \pgfpathmoveto{\pgfpoint{-1.5cm}{-0.75cm}}\pgfpathlineto{\pgfpoint{-1.1cm}{-0.75cm}}
  }
}
\makeatother
\begin{document}
\begin{circuitikz}[american, line width=0.8pt]
\ctikzset{bipoles/length=1.2cm}
\ctikzset{grounds/scale=1.36}
\coordinate (b1) at (0,-2);
\coordinate (e1) at (0,-8);
\coordinate (b4) at (6,-2);
\coordinate (b5) at (8,-2);
\coordinate (g1) at (0,-12);
\coordinate (g3) at (4,-12);
\coordinate (g5) at (8,-12);
\coordinate (j1) at (0,-18);
\coordinate (j4) at (6,-18);
\coordinate (c10) at (18,-4);
\coordinate (d5) at (8,-6);
\coordinate (h6) at (10,-14);
\coordinate (f7) at (12,-10);
\coordinate (g8) at (14,-12);
\draw (b1) to[esource, l_=$V_{1}$, a^=$4\mbox{.}5\,\mathrm{V}$] (e1); % line 3
\draw (-0.081,-4.865) -- ++(0.162,0); % line 3
\draw (0,-4.946) -- ++(0,0.162); % line 3
\draw (-0.081,-5.135) -- ++(0.162,0); % line 3
\node[ground] at (e1) {}; % line 4
\draw (b1) to[sD, l_=$D_{2}$, a^=$\mathrm{1N5817}$] (b4); % line 5
\draw (b5) node[ocirc]{} node[above left]{VSYS}; % line 6
\draw (g1) node[ocirc]{} node[above left]{GP15}; % line 7
\draw (g1) to[R, l_=$R_{1}$, a^=$330\,\Omega$] (g3); % line 8
\draw (g3) to[leD, l_=$D_{1}$] (g5); % line 9
\node[anchor=south] at (6,-11.2) {$\mathrm{red}$}; % line 9
\node[ground] at (g5) {}; % line 10
\draw (j1) node[ocirc]{} node[above left]{GP14}; % line 11
\draw (j1) to[nopb, l_=$S_{W1}$] (j4); % line 12
\node[ground] at (j4) {}; % line 13
\node[dev4w11n8, draw, font=\scriptsize] (part-U1) at (c10) {$\mathrm{Pico\ 2}$}; % line 14
\node[anchor=south] at (part-U1.north) {$U_{1}$}; % line 14
\node[circuitnotemark, font=\scriptsize, anchor=west, xshift=2pt, inner sep=0] at (part-U1.bpin 1) {X}; % line 14
\node[circuitnotemark, font=\scriptsize, anchor=west, xshift=2pt, inner sep=0] at (part-U1.bpin 2) {X}; % line 14
\node[circuitnotemark, font=\scriptsize, anchor=west, xshift=2pt, inner sep=0] at (part-U1.bpin 3) {X}; % line 14
\node[circuitnotemark, font=\scriptsize, anchor=west, xshift=2pt, inner sep=0] at (part-U1.bpin 4) {X}; % line 14
\draw (d5) node[ocirc]{} node[above left]{VSYS}; % line 19
\draw (h6) node[ocirc]{} node[above left]{GP15}; % line 20
\draw (f7) node[ocirc]{} node[above left]{GP14}; % line 21
\node[ground] at (g8) {}; % line 22
\draw (b4) -- (b5); % line 24
\draw (part-U1.pin 1) -| (d5); % line 25
\draw (part-U1.pin 2) -| (h6); % line 26
\draw (part-U1.pin 3) -| (f7); % line 27
\draw (part-U1.pin 4) -| (g8); % line 28
\node[anchor=south west, inner sep=2pt, yshift=4pt, circuitnotemark, font=\LARGE\bfseries] (circuittitle) at (current bounding box.north west) {X};
\path (circuittitle.south west) rectangle ++(15.21,0.853);
\node[anchor=north east, inner sep=2pt, circuitnotemark, font=\large] (circuitstamp) at ([yshift=-0.593cm]current bounding box.south east) {X};
\path (circuitstamp.north east) rectangle ++(-5.08,-0.593);
\end{circuitikz}
\end{document}
